% ------------------------ CONCLUSION ---------------------------------
\section{Conclusion}
\label{sec:conc}

We have shown that semi-annual changes to OSEBX are predictable out-of-sample at horizons up to 100 trading days before the effective date, using only the four constituent-selection variables that Euronext itself publishes \parencite{OSEBX-rule-book}. A logistic GLM \parencite{glm-generalized-linear-models} and an XGBoost classifier \parencite{xgboost-a-scaleable-tree-boosting-system} both exceed 92\% accuracy and 0.97 AUROC, and equally-weighted long--short portfolios that buy predicted additions and short predicted deletions deliver 0.95\% (GLM) and 0.32\% (XGBC) monthly excess return over OSEBX from 2010 to 2022, with Fama--French alphas significant at the 5\% level. Embedded as the active sleeve of a 2\%-tracking-error enhanced-index portfolio, the XGBC variant retains a 10\%-significant alpha; the higher-return but lower-capacity GLM variants do not.

The methodological contribution that drives the result is the conditional posterior probability threshold of equation~(\ref{eq:cppt}). It formalises the fact that, in a highly persistent classification problem, the cost of a missed signal is structurally larger than the cost of a false alarm, and it gives the analyst one interpretable knob, $\varepsilon$, with which to trade calibration accuracy for signal density. The empirical knife-edge --- the 95\%-accurate baseline is unprofitable, the 92\%-accurate XGBC variant clears every risk-adjustment test we run --- is a useful caution against optimising machine-learning classifiers on accuracy alone when the downstream decision is a portfolio.

Three limitations bound the result. The Norwegian factors of \textcite{bernt-fama-french} are constructed on the wider Oslo Stock Exchange and only approximate the relevant factor structure. The 4 bps transaction-cost assumption \parencite{dnb-transaction-costs} is institutional and would understate retail costs by a factor of five or more. And the index effect itself may attenuate as more participants exploit it \parencite{spglobal-what-happened-to-the-index-effect}; replicating the analysis on the next five rebalancing events is the natural robustness check. Subject to those caveats, the strategy is implementable today, with publicly available rule-book data and an open-source toolchain.
